% =====================================================================
\chapter{Los números reales}
\label{tema:1}
% =====================================================================

Este tema fija las reglas de juego de todo el módulo: qué se puede hacer con los números reales y por qué. En lugar de dar por sabidas las propiedades de la aritmética, se parte de una lista corta de axiomas, la de un \emph{cuerpo ordenado}, y se deducen de ella las reglas habituales. La presentación sigue los capítulos~1 y~2 de Spivak~\cite{spivak}; las secciones de potencias y radicales siguen los §1.1--1.3 de OpenStax~\cite{openstax}.

\section{Conjuntos numéricos}

Se trabaja con la cadena de inclusiones $\N\subset\Z\subset\Q\subset\R$, donde $\N=\{1,2,3,\dots\}$ son los naturales, $\Z$ los enteros, $\Q=\{p/q: p,q\in\Z,\ q\neq 0\}$ los racionales y $\R$ los reales. Un real que no es racional se llama \emph{irracional}.

\begin{proposicion}\label{prop:decimal-periodico}
Un número real es racional si y solo si su expresión decimal es finita o periódica.
\end{proposicion}

\begin{proof}[Idea de la demostración]
Si $x=p/q$ con $q>0$, la división larga de $p$ entre $q$ solo puede producir los restos $0,1,\dots,q-1$; en a lo sumo $q$ pasos un resto se repite y, desde ese momento, las cifras se repiten. Recíprocamente, si $x=0{,}\overline{d_1\dots d_k}$, entonces $10^{k}x-x=(d_1\dots d_k)$ es un entero, luego $x=(d_1\dots d_k)/(10^{k}-1)\in\Q$; el caso de un anteperiodo se reduce a este multiplicando por una potencia de~10.
\end{proof}

\begin{ejemplo}\label{ej:periodico}
Sea $x=0{,}\overline{36}$. Entonces $100x=36{,}\overline{36}$ y $100x-x=36$, de donde $x=36/99=4/11$.
\end{ejemplo}

\section{Axiomas de cuerpo}

\begin{definicion}[Cuerpo]\label{def:cuerpo}
Un \emph{cuerpo} es un conjunto $K$ con dos operaciones, suma $(a,b)\mapsto a+b$ y producto $(a,b)\mapsto ab$, que cumplen, para todos $a,b,c\in K$:
\begin{enumerate}[label=(C\arabic*)]
  \item asociatividad: $(a+b)+c=a+(b+c)$ y $(ab)c=a(bc)$;
  \item conmutatividad: $a+b=b+a$ y $ab=ba$;
  \item elementos neutros: existen $0,1\in K$, con $0\neq 1$, tales que $a+0=a$ y $a\cdot 1=a$;
  \item opuesto: para cada $a$ existe $-a$ con $a+(-a)=0$;
  \item inverso: para cada $a\neq 0$ existe $a^{-1}$ con $aa^{-1}=1$;
  \item distributividad: $a(b+c)=ab+ac$.
\end{enumerate}
Se definen la resta $a-b=a+(-b)$ y, si $b\neq 0$, la división $a/b=ab^{-1}$.
\end{definicion}

$\Q$ y $\R$ son cuerpos; $\Z$ no lo es, porque $2$ no tiene inverso en $\Z$. Todas las reglas de la aritmética escolar se deducen de (C1)--(C6).

\begin{proposicion}[Consecuencias de los axiomas de cuerpo]\label{prop:cuerpo}
Para todos $a,b,x$ de un cuerpo:
\begin{enumerate}
  \item si $a+x=a$, entonces $x=0$ (unicidad del neutro); si $a+x=0$, entonces $x=-a$ (unicidad del opuesto);
  \item $a\cdot 0=0$;
  \item $(-a)b=-(ab)$ y $(-a)(-b)=ab$;
  \item si $ab=0$, entonces $a=0$ o $b=0$.
\end{enumerate}
\end{proposicion}

\begin{proof}
(i) Sumando $-a$ a ambos lados de $a+x=a$ y usando (C1), (C4) y (C3): $x=0+x=(-a+a)+x=-a+(a+x)=-a+a=0$. El segundo apartado es análogo: $x=(-a+a)+x=-a+(a+x)=-a+0=-a$.

(ii) Por (C3) y (C6), $a\cdot 0=a(0+0)=a\cdot 0+a\cdot 0$; por (i) aplicado a $a\cdot 0$, se obtiene $a\cdot 0=0$.

(iii) $ab+(-a)b=(a+(-a))b=0\cdot b=0$, luego, por la unicidad del opuesto, $(-a)b=-(ab)$. Aplicándolo dos veces, $(-a)(-b)=-(a(-b))=-(-(ab))=ab$, donde $-(-c)=c$ porque $c+(-c)=0$ dice que $c$ es el opuesto de $-c$.

(iv) Si $a\neq 0$, se multiplica $ab=0$ por $a^{-1}$: $b=(a^{-1}a)b=a^{-1}(ab)=a^{-1}\cdot 0=0$.
\end{proof}

El apartado~(iv) es el que justifica resolver $(x-2)(x+3)=0$ igualando a cero cada factor.

\section{Axiomas de orden}

\begin{definicion}[Cuerpo ordenado]\label{def:orden}
Un cuerpo $K$ es \emph{ordenado} si existe un subconjunto $P\subset K$, el de los \emph{positivos}, tal que:
\begin{enumerate}[label=(O\arabic*)]
  \item tricotomía: para cada $a\in K$ se cumple exactamente una de $a=0$, $a\in P$, $-a\in P$;
  \item si $a,b\in P$, entonces $a+b\in P$;
  \item si $a,b\in P$, entonces $ab\in P$.
\end{enumerate}
Se escribe $a<b$ (o $b>a$) si $b-a\in P$, y $a\le b$ si $a<b$ o $a=b$.
\end{definicion}

\begin{proposicion}[Reglas de las desigualdades]\label{prop:orden}
En un cuerpo ordenado, para todos $a,b,c$:
\begin{enumerate}
  \item si $a<b$ y $b<c$, entonces $a<c$;
  \item si $a<b$, entonces $a+c<b+c$;
  \item si $a<b$ y $c>0$, entonces $ac<bc$;
  \item si $a\neq 0$, entonces $a^2>0$; en particular, $1>0$.
\end{enumerate}
\end{proposicion}

\begin{proof}
(i) $c-a=(c-b)+(b-a)$ es suma de dos positivos, luego positivo por (O2). (ii) $(b+c)-(a+c)=b-a\in P$. (iii) $bc-ac=(b-a)c$ es producto de dos positivos, positivo por (O3). (iv) Por (O1), $a\in P$ o $-a\in P$. En el primer caso $a^2=a\cdot a\in P$; en el segundo, $a^2=(-a)(-a)\in P$ por la \cref{prop:cuerpo}(iii). Como $1=1^2$ y $1\neq 0$, se tiene $1>0$.
\end{proof}

\begin{observacion}[La completitud]
$\Q$ y $\R$ son ambos cuerpos ordenados; lo que los distingue es el \emph{axioma del supremo}: todo subconjunto no vacío de $\R$ acotado superiormente tiene una cota superior mínima~\cite[cap.~8]{spivak}. Gracias a él existen en $\R$ números como $\sqrt 2$, que no están en $\Q$ (\cref{teo:raiz2}). En este módulo el axioma solo se usa implícitamente, al admitir que existen las raíces $n$-ésimas y los logaritmos de los reales positivos.
\end{observacion}

\section{Números irracionales}

\begin{lema}\label{lem:par}
Sea $n\in\Z$. Si $n^2$ es par, entonces $n$ es par.
\end{lema}

\begin{proof}
Por contrarrecíproco. Si $n$ es impar, $n=2k+1$ con $k\in\Z$, y $n^2=4k^2+4k+1=2(2k^2+2k)+1$ es impar.
\end{proof}

\begin{teorema}\label{teo:raiz2}
$\sqrt2$ es irracional.
\end{teorema}

\begin{proof}
Por reducción al absurdo, supóngase $\sqrt2=p/q$ con $p,q\in\Z$, $q\neq0$, y la fracción irreducible. Elevando al cuadrado, $p^2=2q^2$, luego $p^2$ es par y, por el \cref{lem:par}, $p=2r$. Sustituyendo, $4r^2=2q^2$, es decir, $q^2=2r^2$, y por el mismo lema $q$ es par. Entonces $p$ y $q$ son ambos pares, en contra de que la fracción sea irreducible.
\end{proof}

\section{Potencias, radicales e intervalos}

Para $a\in\R$ y $n\in\N$ se define $a^n=a\cdots a$ ($n$ factores), $a^0=1$ si $a\neq 0$, y $a^{-n}=1/a^n$. Para $a>0$ se define $a^{1/n}=\sqrt[n]{a}$ como el único real positivo cuya potencia $n$-ésima es $a$, y $a^{m/n}=(a^{1/n})^m$.

\begin{proposicion}[Leyes de los exponentes]\label{prop:exponentes}
Para $a,b>0$ y $r,s\in\Q$: $a^{r}a^{s}=a^{r+s}$, $(a^{r})^{s}=a^{rs}$, $(ab)^{r}=a^{r}b^{r}$ y $a^{-r}=1/a^{r}$. Para exponentes enteros, las tres primeras valen también con $a,b\neq 0$.
\end{proposicion}

Para exponentes naturales, las leyes se prueban por inducción sobre el exponente; para exponentes racionales se reducen a ese caso usando la unicidad de la raíz $n$-ésima~\cite[§1.2--1.3]{openstax}.

\begin{definicion}[Intervalos]
Para $a<b$: $[a,b]=\{x: a\le x\le b\}$, $(a,b)=\{x: a<x<b\}$, y análogamente $[a,b)$, $(a,b]$, $(a,\infty)$, $(-\infty,b]$, etc. Los símbolos $\pm\infty$ no son números reales y nunca llevan corchete cerrado.
\end{definicion}

\begin{observacion}[Errores frecuentes]
\begin{enumerate}[label=(\alph*)]
\item $\sqrt{a^2}=\abs{a}$, no $a$: con $a=-3$, $\sqrt{9}=3$.
\item $(a+b)^2\neq a^2+b^2$ y $\sqrt{a+b}\neq\sqrt a+\sqrt b$ en general.
\item Al multiplicar una desigualdad por un número negativo, su sentido se invierte (\cref{eje:1-negativo}).
\item No se puede dividir por una expresión que puede valer cero: de $x^2=x$ no se sigue $x=1$, porque se pierde la solución $x=0$.
\end{enumerate}
\end{observacion}

\begin{observacion}[Conexión con el curso]
Los ordenadores no trabajan con $\R$ sino con un conjunto finito de números en coma flotante (\cref{tema:6}). En ese conjunto la suma no es asociativa: con el formato bfloat16, $(256+1)+1=256$ mientras que $256+(1+1)=258$ (\cref{eje:6-espaciado}). Por eso el resultado de sumar los gradientes de muchas GPU depende del orden de la reducción (módulos~11 y~17), y por eso los axiomas de esta sección son una idealización que hay que vigilar en la práctica.
\end{observacion}

\begin{lecturas}
\item Spivak~\cite{spivak}, cap.~1 «Basic Properties of Numbers» (axiomas P1--P12 y sus consecuencias) y cap.~2 «Numbers of Various Sorts» (irracionalidad de $\sqrt2$).
\item OpenStax \emph{College Algebra}~\cite{openstax}, §1.1 «Real Numbers: Algebra Essentials», §1.2 «Exponents and Scientific Notation» y §1.3 «Radicals and Rational Exponents».
\item Hammack~\cite{hammack}, cap.~6 «Proof by Contradiction», para practicar demostraciones de irracionalidad.
\end{lecturas}

\section*{Ejercicios del tema 1}
\addcontentsline{toc}{section}{Ejercicios del tema 1}

\begin{ejercicio}{1}{eje:1-clasificar}
Indique a cuáles de los conjuntos $\N$, $\Z$, $\Q$ y $\R\setminus\Q$ pertenece cada número: $-3$, $0{,}75$, $0{,}\overline{3}$, $\sqrt9$, $\sqrt8$, $\pi$, $1/7$ y $0$.
\end{ejercicio}

\begin{ejercicio}{1}{eje:1-axiomas}
En la siguiente deducción de «si $a+x=a+y$, entonces $x=y$», indique el axioma (C1)--(C6) que justifica cada igualdad:
\[
x\overset{(1)}{=}0+x\overset{(2)}{=}((-a)+a)+x\overset{(3)}{=}(-a)+(a+x)\overset{(4)}{=}(-a)+(a+y)\overset{(5)}{=}((-a)+a)+y\overset{(6)}{=}0+y\overset{(7)}{=}y.
\]
\end{ejercicio}

\begin{ejercicio}{1}{eje:1-vf}
Razone si son verdaderas o falsas: (a) la suma de dos irracionales es irracional; (b) el producto de un racional no nulo por un irracional es irracional; (c) entre dos racionales distintos siempre hay otro racional.
\end{ejercicio}

\begin{ejercicio}{2}{eje:1-fracciones}
Escriba como fracción irreducible: (a) $0{,}1\overline{6}$; (b) $2{,}\overline{45}$; (c) $0{,}\overline{142857}$.
\end{ejercicio}

\begin{ejercicio}{2}{eje:1-potencias}
Simplifique sin calculadora: (a) $\dfrac{2^{-3}\cdot 4^{2}}{8^{-1}}$; (b) $27^{2/3}\cdot 16^{-3/4}$; (c) $\sqrt{50}-\sqrt{18}+\sqrt{8}$.
\fuente{adaptado de OpenStax \emph{College Algebra}~\cite{openstax}, §1.2 y §1.3.}
\end{ejercicio}

\begin{ejercicio}{2}{eje:1-racionalizar}
Racionalice los denominadores: (a) $\dfrac{1}{\sqrt3-\sqrt2}$; (b) $\dfrac{2}{1+\sqrt5}$; (c) $\dfrac{6}{\sqrt[3]{2}}$.
\fuente{adaptado de OpenStax \emph{College Algebra}~\cite{openstax}, §1.3.}
\end{ejercicio}

\begin{ejercicio}{2}{eje:1-ordenar}
Ordene de menor a mayor $\sqrt2+\sqrt3$, $\sqrt{10}$ y $\pi$ sin calculadora, usando solo que $3{,}1415<\pi<3{,}1416$. (Indicación: compare los cuadrados y aplique el \cref{eje:1-potencia-n}.)
\end{ejercicio}

\begin{ejercicio}{2}{eje:1-cientifica}
(a) Calcule $\dfrac{(3{,}2\cdot 10^{5})(4\cdot 10^{-8})}{1{,}6\cdot 10^{-2}}$ en notación científica. (b) Un modelo tiene $7\cdot 10^{10}$ parámetros y cada uno ocupa 2 bytes. ¿Cuántos bytes ocupa? Expréselo en gigabytes ($1\ \text{GB}=10^{9}$ bytes).
\fuente{adaptado de OpenStax \emph{College Algebra}~\cite{openstax}, §1.2 (notación científica).}
\end{ejercicio}

\begin{ejercicio}{2}{eje:1-intervalos}
Sean $A=[-2,3)$ y $B=(1,5]$. Escriba como intervalos o uniones de intervalos: $A\cap B$, $A\cup B$, $A\setminus B$ y $\R\setminus A$.
\end{ejercicio}

\begin{ejercicio}{2}{eje:1-ecuacion-producto}
Resuelva $x^3=x$ en $\R$ justificando cada paso con la \cref{prop:cuerpo}. Explique qué error se comete al «simplificar» dividiendo por $x$.
\end{ejercicio}

\begin{ejercicio}{3}{eje:1-quebrados}
Demuestre, a partir de los axiomas de cuerpo, que si $b\neq0$ y $d\neq0$, entonces
\[
\frac ab+\frac cd=\frac{ad+bc}{bd}\qquad\text{y}\qquad \frac ab\cdot\frac cd=\frac{ac}{bd}.
\]
(Puede usar sin demostración que $(bd)^{-1}=b^{-1}d^{-1}$.)
\fuente{adaptado de Spivak~\cite{spivak}, cap.~1, problemas finales.}
\end{ejercicio}

\begin{ejercicio}{3}{eje:1-negativo}
Demuestre, a partir de los axiomas de orden: (a) si $a<b$ y $c<0$, entonces $ac>bc$; (b) si $a>0$, entonces $1/a>0$; (c) si $0<a<b$, entonces $1/b<1/a$.
\fuente{adaptado de Spivak~\cite{spivak}, cap.~1, problemas finales.}
\end{ejercicio}

\begin{ejercicio}{3}{eje:1-potencia-n}
Demuestre por inducción que, si $0\le a<b$, entonces $a^n<b^n$ para todo $n\in\N$. Deduzca que, si $a,b\ge 0$ y $a^2<b^2$, entonces $a<b$.
\fuente{adaptado de Spivak~\cite{spivak}, cap.~1, problemas finales.}
\end{ejercicio}

\begin{ejercicio}{3}{eje:1-raiz3}
Demuestre que $\sqrt3$ es irracional y, a partir de ello, que $\sqrt2+\sqrt3$ también lo es.
\fuente{adaptado de Spivak~\cite{spivak}, cap.~2, problemas finales.}
\end{ejercicio}

\begin{ejercicio}{3}{eje:1-epsilon}
Sean $x,y\in\R$. Demuestre: (a) si $a<b$, entonces $a<\frac{a+b}{2}<b$; (b) si $x\le y+\varepsilon$ para todo $\varepsilon>0$, entonces $x\le y$.
\end{ejercicio}
